\documentclass[11pt,a4paper]{article}
% Stage 10 manuscript draft. Frozen theory: V2B-THEORY-FREEZE-2026-10-21-v1.
% This document is NOT submission-ready and is subject to Stage 11--14 review.
\usepackage[margin=1in]{geometry}
\usepackage{amsmath,amssymb,amsthm}
\usepackage[round,authoryear]{natbib}
\usepackage[hidelinks]{hyperref}
\usepackage[T1]{fontenc}
\usepackage{lmodern}
\usepackage{microtype}
\usepackage{booktabs}

\newtheorem{proposition}{Proposition}
\newtheorem{lemma}[proposition]{Lemma}
\newtheorem{corollary}[proposition]{Corollary}
\newcommand{\SW}{\mathrm{SW}}
\newcommand{\SU}{\mathrm{SU}}
\newcommand{\IS}{\mathrm{IS}}
\newcommand{\CS}{\mathrm{CS}}

\title{International Standards Coalitions and Endogenous Product Design}
\author{}
\date{}
\begin{document}
\maketitle

\begin{abstract}
Can a welfare-improving international technical standard fail
to emerge because of private choices made \emph{after} governments
commit to compatibility? We construct a five-stage three-country
standards game in which two countries first choose a technical
compatibility depth, a third may apply to join, and firms subsequently
choose costly nonstandardized product attributes before Cournot
competition. At an exactly specified parameter witness, endogenous
private design makes every strictly positive bilateral union
unattractive to its members and accession unattractive to the
outsider. Nonetheless, a hypothetical three-country agreement
strictly Pareto-dominates fragmentation even after firms
reoptimize their product attributes. A controlled forecast that
freezes product designs at their fragmented-market choices
instead predicts an accepted international agreement. A separately
certified nonlinear design mapping preserves this existential
contrast, although the feasible compatibility depth changes
and a pre-registered same-depth robustness test fails.
The result is an institution-specific, computer-assisted
existence argument, not a universal standardization theorem
or a claim of empirically observed accession refusal.
\end{abstract}

\section{Introduction}\label{sec:introduction}

Can international compatibility fail to emerge even when
each country would gain from a common technical standard?
In industries where interoperability links consumers or
networks across products, a government may expect
compatibility to raise consumer welfare. The same
commitment can change firms' incentives to differentiate
the nonstandardized portions of their products, however,
and those decisions are made \emph{after} governments
announce a technical standard. A forecast that treats
other product characteristics as fixed may therefore
give the wrong answer about which international
agreements countries will accept.

We study an intentionally stylized setting in which
two countries first negotiate a technical-compatibility
union and a third may subsequently seek accession.
The existing members must unanimously approve the
third, and the new entrant inherits the original
compatibility depth. Each country values its domestic
consumers together with the global profits of its
home firm. After the government's decision, firms
independently choose costly residual product attributes,
and then compete in quantities across three separate
national markets. In this environment, an interoperable
\emph{interface} need not make final products identical.

The paper constructs a restricted but sharp
equilibrium counterexample. At one exact rational
parameter vector, with the firms freely reoptimizing
their nonstandardized attributes, no bilateral union
at any strictly positive compatibility depth is
acceptable to its prospective members, and the
third country would rather remain outside than join
the corresponding enlarged international union.
Thus the specified bilateral-first procedure implements
no substantive technical agreement. Yet an
\emph{international} compatible allocation at one
positive depth, conditional on firms subsequently
reoptimizing, would improve \emph{every country's}
welfare relative to complete technical fragmentation.
The obstacle is the path of coalition formation and
the outsider's intermediate option, not a negative
aggregate value of technical compatibility.

A controlled benchmark reveals the role of downstream
design. Hold each firm's residual product attributes
fixed at the design it chooses under separate
standards, while leaving the governments' actions,
market competition and once-paid design costs
unchanged. In that \emph{strategy-restricted}
counterfactual game, a positive-depth international
union is now acceptable to the outsider and the
incumbent members and is preferred by the proposer
to any terminal bilateral union. The predicted
institutional regime therefore reverses when
endogenous product attributes are restored.
Our main all-depth inequalities are supported by
analytical continuation arguments and an independently
reconstructed rational-interval certificate; they
do not rest merely on a numerical policy grid.

The mechanism is not two independently novel
sign-flip results for the incumbent and the outsider.
At the symmetric incompatibility point their marginal
changes are linked by a simple graph-symmetry identity.
Rather, policy-dependent private residual design
changes product-market competition and the national
valuation of exported profits, while the
bilateral-first participation institution maps
the resulting country payoffs into a blocked
multilateral agreement. The government can experience
a positive direct benefit from compatibility yet a
negative \emph{total} derivative once firms adjust
their residual characteristics.

The distinction is important relative to several
prior theories. Network benefits versus heightened
competition and standards-coalition formation are
already familiar \citep{EconomidesFlyer1998,
EconomidesSkrzypacz2003}. Three-country standard
recognition and regional unions were studied by
\citet{GandalShy2001}, with exogenous Salop variety
and subsequent Bertrand pricing.
\citet{CabralSalant2014} show that standardization
can weaken firms' innovation incentives,
and \citet{BuchholzHattori2024} identify a
technology--coalition paradox in public-good
cooperation. \citet{TakaradaKawabata2025}
examine regional versus multilateral regulatory
mutual recognition and the significance of
exogenous differentiation. We do not claim
priority for these general ingredients.
The narrower contribution is to exhibit and
certify a combination of \emph{governmental
sequencing, post-commitment private design,
and Pareto-beneficial but inaccessible final
compatibility} in an explicit market and
national-welfare model.

Our conclusion is correspondingly limited.
The exact model has symmetric countries,
linear-quadratic conditional demand, and
a specified bilateral-first admission rule.
A separately pre-registered experiment with
a nonlinear mapping from product design to
substitutability reproduces the \emph{existence}
contrast at the same economic parameters,
but not at the same compatibility depth:
its initially specified fixed-depth participation
condition fails. No generic result for all
product-design functions or bargaining
institutions is established. Historical
bilateral technical initiatives such as
GSM development and ChaoJi charging cooperation
motivate the timing of early technical alignment
before wider standardization; they are not
claimed to be empirical realizations of the
model's third-country rejection.

Section~\ref{sec:model} defines the
sequential game and separates standardized
interfaces from residual product attributes.
Section~\ref{sec:equilibrium} derives the
Cournot and design continuations,
Section~\ref{sec:results} establishes
the endogenous and frozen-design
coalition results, and Section~\ref{sec:welfare}
shows the remaining Pareto improvement.
Section~\ref{sec:robustness} gives the
nonlinear diagnostic, followed by the
closest-literature comparison, institutional
limitations and proof appendix.

\section{Model}\label{sec:model}

\subsection{A bilateral-first technical agreement}

Consider three countries, indexed by $i\in\{A,B,C\}$, each hosting a
single firm, and three separated national product markets. All three
firms sell in each national market. The governments observe all previous
decisions and maximize national welfare. At the technology-policy stage,
compatibility among \emph{interfaces} is distinguished from the
similarity of the other attributes chosen by firms.

The extensive form has five stages. First, government $A$ offers a
bilateral standards union between $A$ and $B$, denoted $\SU$, at
a \emph{technical compatibility depth} $g\in[0,1]$. Government $B$
accepts or rejects; rejection leaves all countries with separate
standards ($\SW$). Second, after an accepted bilateral union, country
$C$ decides whether to request accession. Nonapplication leaves $\SU$
in place. Third, if $C$ applies, both $A$ and $B$ must approve.
Unanimity produces an international standards union ($\IS$), while
either veto retains $\SU$. Accession \emph{inherits the existing $g$}:
no party can renegotiate its depth at this stage.

Fourth, given the realized institutional regime, each firm simultaneously
chooses a residual product-homogeneity index $x_i\in[0,1]$. A higher
$x_i$ reduces differentiation in nonstandardized attributes; critically,
there is no restriction $x_i\geq g$ and no requirement that whole
products be identical merely because their technical interfaces work
together. A firm pays the real design cost
\begin{equation}\label{eq:designcost}
 K_i(x_i)=\frac{\kappa}{2}(1-x_i)^2,\qquad \kappa>0,
\end{equation}
\emph{once across its worldwide operations}. Finally, firms
simultaneously select nonnegative Cournot quantities in every market,
taking expected network use as fixed when choosing quantities.
All firms and governments know the primitives, chosen designs, and
compatibility commitments.

Government votes, applications, and acceptance at indifference are
weakly affirmative. Membership labels at $g=0$ may therefore depend
on indifferent choices even though the corresponding allocations are
identical. We study the \emph{substantive} formation of unions with
$g>0$, not the uniqueness of zero-depth institutional labels.
There are no side payments, adjustment fees, border barriers,
government mandates, or renegotiation after $C$ applies.

\subsection{Demand and distinct dimensions of compatibility}

In every national market, a representative quasilinear consumer has
the conditional utility
\begin{equation}\label{eq:utility}
 U(q\mid n^e)
 =a\mathbf{1}^{\top}q-\frac{1}{2}q^{\top}Bq
       +\eta q^{\top}H n^e,\qquad
 B=\rho H+(1-\rho)X ,
\end{equation}
where $q=(q_A,q_B,q_C)^\top$ is the local consumption vector,
$n^e$ is the \emph{expected} vector of network use in the same
market, $0<\rho<1$ measures the importance of technical
compatibility for substitutability, and $0\leq\eta<1$ measures
network benefits. Throughout, own diagonal entries are
$H_{ii}=X_{ii}=1$. For $i\ne j$, define
\begin{equation}\label{eq:matrices}
 H_{ij}=
 \begin{cases}
 g,&\text{$i$ and $j$ belong to the same standards union},\\
 0,&\text{otherwise},
 \end{cases}
 \qquad
 X_{ij}=\frac{x_i+x_j}{2}.
\end{equation}
Thus $\SW$ has no off-diagonal entries of $H$, $\SU$ adds only the
$AB$ edge, and $\IS$ gives every pair the same inherited $g$.
Residual similarity $X$ need not vanish under separate standards,
and the common interface $H$ need not eliminate other product
differences. The quadratic matrices are positive semidefinite
at the certified parameter witness and design domain.

Holding expectations $n^e$ fixed for the consumer's conditional
choice, inverse demand is
\begin{equation}\label{eq:demand}
 p=a\mathbf{1}-Bq+\eta H n^e.
\end{equation}
Firms share constant marginal production cost $c$. In the Cournot
stage, an individual firm's quantity deviation does \emph{not}
change $n^e$ contemporaneously; equilibrium expectations are
subsequently fulfilled, $n^e=q$. All national markets have
identical primitives and hence identical fulfilled-expectations
quantities in the unique continuation considered below.

\subsection{Profit, welfare, and benchmarks}

Let $d=a-c>0$. At an interior fulfilled-expectations Cournot
equilibrium the margin in each national market is $p_i-c=q_i$.
Consequently, firm $i$ receives operating profits $3q_i^2$
from the \emph{three} markets but pays its design expenditure only
once. Its Stage-4 objective is
\begin{equation}\label{eq:pi}
 \Pi_i(x;S,g)=3q_i(x;S,g)^2-K_i(x_i).
\end{equation}
Conditional consumer surplus in \emph{one} country's market is
$\CS(q;S,g,x)=\tfrac12q^\top Bq$, evaluated at fulfilled
expectations; adding $\eta q^\top Hq$ again would double-count
network utility in this surplus calculation. Government $i$
values its own consumers and its home firm's worldwide net income:
\begin{equation}\label{eq:welfare}
 W_i(S,g;x)=
 \underbrace{\frac12q^\top Bq}_{\text{domestic consumer surplus}}
 +\underbrace{3q_i^2}_{\text{home firm's world operating profit}}
 -\underbrace{\frac{\kappa}{2}(1-x_i)^2}_{\text{once-paid design cost}}.
\end{equation}
World welfare is
\begin{equation}\label{eq:world}
 {\cal W}=3\CS+3\sum_i q_i^2
                       -\frac{\kappa}{2}\sum_i(1-x_i)^2.
\end{equation}

The actual game allows post-policy design responses at every
history. We also analyze a \emph{different restricted-strategy
counterfactual game} in which each $x_i$ is forced to equal the
corresponding $\SW$ design Nash choice after \emph{every} history.
This is a forecast with frozen private designs, not another
equilibrium of the original unconstrained design game. Separately,
we evaluate the conditional welfare of an $\IS$ allocation at a
specified $g$ even if the governments' sequential consent rules
prevent its implementation. That benchmark is not a planner's
unrestricted first best.

Our result is an existence result at the following exact rational
parameter vector:
\begin{equation}\label{eq:witness}
 \theta_0=(d,\eta,\rho,\kappa)
 =\left(1,\frac{379}{500},\frac9{10},\frac{1701}{2000}\right).
\end{equation}
No universal claim is made for arbitrary demand or parameters.

\section{Market and Design Equilibria}\label{sec:equilibrium}
\emph{Stage 10 section reserved. Derive fulfilled-expectations Cournot
first-order conditions and design Nash from the frozen primitives;
cite the Stage 4A full-strategy certificate with its exact scope.}

\section{Sequential Compatibility and Coalition Formation}\label{sec:results}
\emph{Stage 10 section reserved. State only the exact-witness,
institution-specific existence results. Full-government continuation and
all-$g$ inequalities are computer-assisted, not full Lean theorems.}

\section{Welfare and the Participation Wedge}\label{sec:welfare}
% Table is generated from the independent frozen V2b evaluator by make table.
\input{generated/welfare_table.tex}
\emph{Stage 10 welfare interpretation and its derivation are not yet drafted.
The table concerns conditional continuation regimes, not agreements
selected by the government game.}

\section{A Nonlinear Design-Function Check}\label{sec:robustness}

The existence claim is not intended to hold for every way
in which nonstandardized attributes enter substitutability.
To probe dependence on the additive form in
\eqref{eq:matrices}, we examine one
\emph{pre-registered} nonlinear design mapping, holding
the government protocol, exact parameter vector
\eqref{eq:witness}, design-cost function \eqref{eq:designcost},
network mechanism, and Cournot stage \emph{unchanged}:
\begin{equation}\label{eq:quadratic-map}
 X^{(2)}_{ii}=1,\qquad
 X^{(2)}_{ij}=\frac{x_i^2+x_j^2}{2}\quad (i\ne j).
\end{equation}
This is not a mere change of variables: the design cost
continues to be quadratic in the \emph{original} $x_i$.
The marginal effect of private design on product substitution
is thus different from that in the baseline.

A separate three-firm pseudo-gradient and equilibrium
calculation is required. The proof checks strict diagonal
concavity in the entire asymmetric design domain and
uniform member/outside equilibrium root inclusion across
all $g\in[0,1]$, rather than substituting $x_i^2$
into a previously selected baseline equilibrium.

\begin{proposition}[Existential persistence under one nonlinear mapping]
\label{prop:quadratic}
At $\theta_0$ and under \eqref{eq:quadratic-map},
the fully endogenous game still satisfies
\[
 m^{(2)}(g)<s^{(2)},\qquad
 i^{(2)}(g)<o^{(2)}(g)\quad\text{for all }g\in(0,1].
\]
However, in the design-frozen counterfactual at the
\emph{positive but smaller} depth $g=1/200$, $\IS$
is strictly accepted and dominates every terminal
$\SU$ and $\SW$ alternative. At $g=1/20$,
an exogenously constituted international agreement
remains Pareto superior to fragmentation even when
firms reoptimize their nonlinear designs.
\end{proposition}

\begin{proof}
A separate rational-interval Rosen/Gershgorin
certificate bounds the symmetric Jacobian of the
actual three-player firm pseudo-gradient on
$[0,1]^3\times[0,1]$ for each institutional regime,
establishing unique Nash design continuations.
Uniform signed root faces and interval implicit
differentiation then verify both strict country
inequalities across closed rational policy intervals.
In the frozen game, the unique $\SW$ design is
approximately $0.610737480$. At $g=1/200$,
the interval-certified gains satisfy
\[
 \widehat{i}^{(2)}(1/200)-
 \widehat{o}^{(2)}(1/200)>9.09\times10^{-7}>0,
\]
and the proposed $\IS$ beats the $\SW$ allocation
and a separately certified full-$g$ upper envelope
for $\SU$ member payoffs. At $g=1/20$ with
\emph{endogenous} designs, the interval-certified
national welfare improvement of the hypothetical
$\IS$ over $\SW$ is approximately $0.016601147>0$.
The separate certificate and reproduced program
targets are described in
Appendix~\ref{app:proofs}.
\end{proof}

\paragraph{A pre-registered failed comparison.}
The initially specified \emph{same-depth} test did
\emph{not} survive the nonlinear mapping: holding
firms' designs frozen and keeping $g=1/20$,
country $C$ would \emph{reject} international
accession, with
\begin{equation}\label{eq:failed-comparison}
 \widehat{i}^{(2)}(1/20)-
 \widehat{o}^{(2)}(1/20)\simeq-0.000440395<0.
\end{equation}
This negative finding must not be reclassified as
passing merely because a \emph{distinct existential}
positive-depth result at $g=1/200$ was subsequently
proved. The latter participation gain is only about
$9.1\times10^{-7}$, so the effective compatibility
range may be quantitatively fragile.

Proposition~\ref{prop:quadratic} documents exactly
\emph{two} certified design maps, $\phi(x)=x$
and $\phi(x)=x^2$. It does not establish a
uniform theorem for $\phi(x)=x^p$ with arbitrary
$p\in[1,2]$, for general nonlinear network adoption,
asymmetric national markets, or alternative bargaining
institutions. Such stronger claims would require
fresh primitives, equilibrium proofs, and separate
certification. Likewise, no numerical parameter
radius for the strictly local open neighborhood of
\eqref{eq:witness} is claimed.

\section{Relation to Existing Theory}\label{sec:literature}

\paragraph{Government standards and outsider incentives.}
\citet{Petropoulou2013} studies national
minimum-quality standards and the inability
in some circumstances to implement a
world-optimal standard by bilateral
adjustments without transfers.
\citet{Costinot2008} compares incomplete
institutional agreements on product
standards under national treatment
versus mutual recognition. Hence neither
governmental standard-setting failure nor
the existence of nonstandardized product
dimensions is an independent priority claim.

The closest governmental predecessors already establish that
regional standards may create incentives incompatible with wider
coordination. \citet{GandalShy2001} study recognition and
regional standards unions in three countries, with government
welfare combining domestic consumers and global profits of the
home firm. Their differentiated brands are fixed on a Salop
circle and compete in prices. Moreover, the absence of a
\emph{formal} union can coexist with worldwide mutual
recognition; their no-union regime is not the no-interface-link
$\SW$ benchmark in this paper.
\citet{TakaradaKawabata2025} model three-country
regional versus multilateral regulatory mutual recognition.
Under independently chosen national standards, regional
coordination can be unattractive, while outsiders can block
multilateral enlargement in other regimes. They also
consider exogenous product differentiation and standards
compliance costs. These are substantive prior results, not
merely similar institutional labels. Our additional margin
is a post-agreement Nash adjustment of \emph{residual} product
attributes that can reverse the country consent margins
under a fixed-depth, bilateral-first agenda; it is not the
discovery that an outsider may prefer to veto accession.

\paragraph{Endogenous differentiation and strategic trade policy.}
\citet{BranderSpencer2015} analyze endogenous
horizontal differentiation in international
Bertrand and Cournot oligopoly, including its
effect on patterns and gains from trade.
\citet{ZhouSpencerVertinsky2002} examine
strategic trade policy with endogenous export
quality, and \citet{JeanneretVerdier1996}
study trade-policy commitment and standards
in vertically differentiated markets.
Policy-dependent private product characteristics
and national exporter-profit effects are
therefore known mechanisms. The present
conditional response and incidence propositions
target the separate question of how the
\emph{same} staged government standardization
institution changes its country-specific
participation feasibility when firms
re-optimize nonstandardized horizontal attributes.

\paragraph{Firm compatibility, quality, and investment.}
The opposition between network benefits and product-market
competition is well established \citep{EconomidesFlyer1998,
EconomidesSkrzypacz2003}.
\citet{CabralSalant2014} establish that standards can reduce
private technology-development incentives. The published
article's proof numbering has not been independently
matched to the accessible earlier working-paper version;
we therefore do not attribute particular numbered
propositions to it.
\citet{BranderSpencer2022} also show that making
horizontal product differentiation endogenous can
reverse familiar Bertrand-versus-Cournot price and
profit rankings. Thus a disagreement between
fixed- and endogenous-differentiation forecasts
is not itself a new generic mechanism; their model
does not study three-country government accession
at inherited technical depth.
More closely, \citet{ChengHuang2025} examine endogenous
product quality and firm veto rights over compatibility in
network industries. Their published abstract reports
incompatibility under strong network effects when firms
hold veto power, together with positive welfare effects
of compatibility. Their accessible publisher appendix
discusses firm-level veto/no-veto subgames.
\citet{BuccellaFantiGori2024} likewise address
compatibility choice and quality differences in Cournot
network goods; a published corrigendum exists, so we
do not rely on an unchecked original theorem beyond
the available scope. Endogenous product characteristics,
private veto, and inefficient incompatibility are
therefore \emph{not} priority claims of the present paper.

\paragraph{Sequence and the conditional forecast.}
\citet{BuchholzHattori2024} examine technology choices
and coalition formation in a public-good game, with the
technology decision preceding coalition formation.
Here the government chooses interface compatibility
first; private firms subsequently optimize an independent,
costly horizontal-design margin; national governments anticipate the resulting domestic surplus
and exporters' rents when making their earlier consent decisions. The
particular result added to the predecessors is the
\emph{controlled forecast reversal} between endogenous
and frozen residual design, together with an inaccessible,
conditional Pareto gain, under the inherited-depth
A/B-first admission procedure.
Backward induction and uniqueness criteria draw on
standard methods \citep{RayVohra1999,Rosen1965}, not
new game-theoretic solution concepts.

This is a \emph{source-scoped} distinction, not a
comprehensive priority certificate. In particular,
the complete main-text propositions of
\citet{ChengHuang2025} were not available for exhaustive
proposition-level matching; this unresolved comparison
is disclosed rather than treated as evidence of novelty.
No result below establishes that all standardization
institutions, firm-design functions, or parameter vectors
exhibit the same failure.

\section{Institutional Interpretation and Limitations}
\label{sec:discussion}

Our extensive form is a stylized description of
\emph{pre-international technical coordination},
not a description of the voting rules of ISO or
IEC. It is plausible for a pair of governments
or nationally backed standards organizations to
coordinate an interface and later consider
broader participation. Historical motivations
include French--German cooperation during the
development of GSM before wider deployment
commitments, and collaboration between the
CHAdeMO Association and the China Electricity
Council on a next-generation charging interface
\citep{GSMAHistory,CHAdeMO2020}.
The latter also illustrates why an interoperable
physical interface need not fix every protocol
or residual product feature. Neither history
is evidence that the exact voting protocol in
Section~\ref{sec:model} operated, or that a
third country historically rejected accession
because of the design mechanism we model.

A positive $g$ in the model represents an
economically binding compatibility decision
that changes cross-product substitution and
network use in the ensuing product market.
Mere information sharing between technical
committees without deployed compatibility
would \emph{not} be sufficient to give a
soft memorandum the meaning of $g>0$.
Actual charging networks have two-sided
adoption and installed-base dynamics, and
actual telecommunications standardization
involves institutions and actors beyond
our three identical countries. These features
are deliberately absent from
\eqref{eq:utility}--\eqref{eq:welfare}.

A useful diagnostic comes from comparing the
government agenda with a distinct hypothetical
procedure. If the three countries were presented
with a \emph{single} unanimously accepted or
rejected $\IS$ offer at $g=.05$, with $\SW$
as the fallback and no intermediate $\SU$,
Proposition~\ref{prop:pareto} establishes that
each country would prefer accepting under
the specified market continuation. The
inaccessibility result is therefore tied to
the \emph{bilateral-first pathway and its
outsider opportunity}; it is not a theorem
that internationally compatible technology
necessarily fails in every bargaining
institution. We have not modelled the
legal authority, commitment mechanism,
transfer scheme, or enforcement that could
make this alternative offer available.
It should not be read as an implementable
policy proposal.

Other boundaries are equally important.
The results are computer-assisted existence
claims at an exact rational witness and
an unquantified local parameter neighborhood,
not a broad calibration to real firms.
Three countries are symmetric; each firm
serves all three identical national markets;
network expectations concern contemporaneous
quantities; demand is quadratic; and the
design game uses a real once-paid
quadratic cost. The certified nonlinear
design check changes only one substitution
mapping, not the demand foundation or
sequential government protocol.
All-g equilibrium/root assertions rely on
rational interval methods, while the separately
built Lean artifact proves four pointwise
payoff comparisons at \emph{rational design
proxies}; it does not formalize the whole
Nash and subgame-perfect equilibrium analysis.

\section{Conclusion}\label{sec:conclusion}

Evaluating a prospective international
technical agreement with firms' product
attributes held fixed can mispredict which
government coalitions are feasible. In a
three-country bilateral-first standards
consent game, we construct and certify a
case in which post-policy private
differentiation rules out every positive
depth of regional or international
agreement, even though a hypothetical
international standard still strictly
improves each country's welfare over
fragmentation. A forecast that fixes
firms' designs at their fragmented-market
choices instead predicts a feasible
international agreement. The qualitative
existence contrast persists under one
separately certified nonlinear residual
design mapping, although the depth of a
feasible frozen-design offer changes
substantially. The contribution is a
specific interaction between private
design responses and a sequential
participation institution, not a
general prescription for international
standardization.

\appendix
\section{Proof and Reproduction Roadmap}\label{app:proofs}
\emph{Stage 10 technical proof manuscript pending. The auditable theorem
certificate mapping is maintained at
\texttt{production/THEOREM\_CERTIFICATE\_MAP.md}; automated certificates
are not a replacement for readable bridge equations and proofs.}


\bibliographystyle{plainnat}
\bibliography{references}
\end{document}
